\Needspace{18\baselineskip}
\subsection{All-obstacle equality on separated products}
\label{subsec:separated-product-all-obstacles}

For sums of curve potentials, equality of the ordinary and ideal envelopes has a coordinatewise criterion even when the obstacle mixes the coordinates. Relative stretching competitors can be added before any obstacle is imposed.

\Needspace{14\baselineskip}
\begin{lemma}\label{lem:product-relative-envelope-classes}\index{Relative envelope class!separated products}
Let $N\geq1$, let $X_j$, $1\leq j\leq N$, be connected compact complex manifolds with smooth real $(1,1)$-forms $\theta_j$, and let $u_j,\varphi_j\in\PSH(X_j,\theta_j)$. On $X=\prod_jX_j$, with projections $\operatorname{pr}_j$, put
\[
 \vartheta=\sum_j\operatorname{pr}_j^*\theta_j,\qquad
 u=\sum_j u_j\circ\operatorname{pr}_j,\qquad
 \varphi=\sum_j\varphi_j\circ\operatorname{pr}_j.
\]
Then
\[
 u\in\mathcal E_{\mathrm{env}}(X,\vartheta,\varphi)
 \quad\Longleftrightarrow\quad
 u_j\in\mathcal E_{\mathrm{env}}(X_j,\theta_j,\varphi_j)
 \quad\text{for every }j.
\]
\end{lemma}
\begin{proof}
Ordinary order of the sums is equivalent to ordinary order in each coordinate. To recover the latter, fix the other coordinates where their potentials are finite and restrict the inequality. Conversely, add the coordinate inequalities.

Suppose each coordinate belongs to its relative class. For $b>1$, choose its stretching competitor $h_j$ as in \cref{def:relative-envelope-class}. The sum $h=\sum_jh_j\circ\operatorname{pr}_j$ is a genuine $\vartheta$-plurisubharmonic function and satisfies
\[
 b^{-1}h+(1-b^{-1})\varphi\leq u.
\]
This convex inequality is meaningful at poles. At $b=1$, use $h=u$. Thus the sum belongs to the relative class.

Conversely, fix $b>1$ and a product competitor $h$. Choose a point where $h$ and all the coordinate potentials are finite; such points have full smooth volume. For the $j$th coordinate, fix the others at that point. The restriction $h^{(j)}$ is genuine and $\theta_j$-plurisubharmonic. Put
\[
 c_j=\sum_{\ell\ne j}u_\ell(x_\ell),\qquad
 d_j=\sum_{\ell\ne j}\varphi_\ell(x_\ell).
\]
Restricting the product inequality gives
\[
 b^{-1}\bigl(h^{(j)}-bc_j+(b-1)d_j\bigr)
       +(1-b^{-1})\varphi_j\leq u_j.
\]
This supplies a coordinate competitor for every $b>1$. Together with ordinary order and the automatic $b=1$ competitor, it proves the converse.
\end{proof}

Keep the curves and smooth admissible backgrounds $\sigma_i$ above, and choose arbitrary $u_i\in\PSH(C_i,\sigma_i)$. Put
\[
 u=\sum_i u_i\circ\pi_i,\qquad
 q_i=P_{\sigma_i}[u_i],\qquad
 r_i=P_{\sigma_i}[u_i]_{\mathcal I},\qquad
 Q=\sum_i r_i\circ\pi_i.
\]
Write $\mu_i$ and $E_i$ for the polar and atomic parts of the positive curvature measure $\sigma_i+\ddc u_i$, respectively.

\Needspace{18\baselineskip}
\begin{proposition}\label{prop:separated-product-all-obstacle-criterion}\index{Envelope!all-obstacle separated product criterion}
For these data, $Q=P_{\gamma_\sigma}[u]_{\mathcal I}$. The following are equivalent:
\[
 \begin{gathered}
 P_{\gamma_\sigma}[u]=P_{\gamma_\sigma}[u]_{\mathcal I},
 \qquad \mu_i=E_i\quad\text{for every }i,\\
 u\in\mathcal E_{\mathrm{env}}(Y,\gamma_\sigma,Q),\\
 P_{\gamma_\sigma,F}[u]=P_{\gamma_\sigma,F}[u]_{\mathcal I}
       \quad\text{for every smooth real }F\text{ on }Y.
 \end{gathered}
\]
The obstacle in the last condition may mix the curve coordinates. Zero curve areas and zero model volume are allowed.
\end{proposition}
\begin{proof}
At zero obstacle, \cref{prop:separated-curve-obstacle-projections} identifies the two envelopes with $\sum_iq_i\circ\pi_i$ and $Q$. Independence of the coordinates and their supremum-zero normalizations identify equality of these sums with $q_i=r_i$ for every $i$. By \cref{cor:curve-arbitrary-background-criterion}, this is equivalent to $\mu_i=E_i$ and to $u_i\in\mathcal E_{\mathrm{env}}(C_i,\sigma_i,r_i)$ for every $i$, including zero-area backgrounds. The preceding lemma identifies these memberships with the product relative class. Finally \cref{prop:ideal-obstacle-reprojection} gives equality at every smooth obstacle, and taking $F=0$ gives the reverse implication. This proves all four equivalences.
\end{proof}

\Needspace{16\baselineskip}
\begin{corollary}\label{cor:auxiliary-product-all-obstacle-criterion}\index{Envelope!nonclosed separated product criterion}
On $M=Y\times C_0\times Z$ in background $\Theta_{\sigma,f}$, with any smooth $f\geq0$, put $U=\Psi+u\circ p$. Then
\[
 \begin{aligned}
 P_{\Theta_{\sigma,f}}[U]&=\Psi+\sum_iq_i\circ\pi_i\circ p,\\
 P_{\Theta_{\sigma,f}}[U]_{\mathcal I}&=\Psi+Q\circ p.
 \end{aligned}
\]
The following are equivalent: equality of these two zero-obstacle envelopes; $\mu_i=E_i$ for every $i$; membership of $U$ in the relative envelope class of $\Psi+Q\circ p$; and equality of its ordinary and ideal envelopes at every smooth real obstacle on $M$. This includes $f=0$ and zero curve areas.
\end{corollary}
\begin{proof}
On $B=Y\times C_0$, with projection $\pi$ to $Y$, use the background $\gamma_{\sigma,f}$ in the proof of \cref{lem:arbitrary-product-model-projections}. Sums of coordinate ordinary competitors give
$P_{\gamma_{\sigma,f}}[u\circ\pi]\geq\sum_iq_i\circ\pi_i\circ\pi$.
For the reverse bound, replace an ordinary competitor by its maximum with a sufficiently low translate of $u\circ\pi$. Its restriction to every $Y\times\{z\}$ is genuine, below zero and ordinarily dominated by $u$. The zero-obstacle separation bounds it by $\sum_iq_i\circ\pi_i$. Unit-divisor reduction in \cref{prop:fibration-obstacle-reduction} proves the upstairs ordinary identity.

Each $q_i$ preserves the ideal class of $u_i$ by \cref{cor:ordinary-preserves-i}. Pullback and addition therefore give $U\sim_{\mathcal I}\Phi_\sigma(q)$. Their ideal envelopes coincide, and \cref{lem:arbitrary-product-model-projections} identifies this envelope with $\Psi+Q\circ p$.

If $\mu_i=E_i$ for every $i$, choose the coordinate stretching competitors from the curve relative classes. Their pullback sum plus $\Psi$ is $\Theta_{\sigma,f}$-plurisubharmonic: its current is the sum of the pulled-back positive curve currents, $(f\circ p)s^*\tau$ and the unit-divisor current. Adding their stretching inequalities gives the one for $U$ relative to $\Psi+Q\circ p$, also at poles; ordinary order follows in the same way. Thus $U$ belongs to that relative class. Ideal obstacle reprojection gives all-obstacle equality. Taking the zero obstacle and using the two proved identities recovers $q_i=r_i$, hence $\mu_i=E_i$. This closes the equivalence.
\end{proof}

These criteria concern sums of curve potentials and their specified product extensions. They do not give coordinatewise formulas for mixed-obstacle envelopes or settle the unrestricted zero-volume implication.
